\textbf{Introducción:} UrbanTracker es una plataforma que optimiza la movilidad urbana mediante el seguimiento en tiempo real de vehículos de transporte público. Integra una aplicación móvil para conductores (desarrollada en React Native) y una plataforma web para usuarios y administradores, comunicados con un backend en Java (Spring Boot) por MQTT. UrbanTracker permite visualizar en un mapa interactivo rutas y vehículos activos, así como gestionar rutas y conductores. En este trabajo se describen los requisitos del sistema, el marco teórico de tecnologías de geolocalización, el enfoque de desarrollo empleado, los resultados preliminares de funcionamiento y las conclusiones. Los resultados indican que la combinación de React Native, Spring Boot y MQTT logra actualizaciones de ubicación con latencias bajas y alta eficiencia, mejorando la experiencia de los usuarios y la gestión operativa del transporte. 
\textbf{Métodos:} Se desarrolló una plataforma compuesta por una aplicación móvil para conductores (con GPS integrado), una aplicación web para usuarios y un panel web administrativo. El backend se implementó en Java Spring Boot exponiendo servicios REST y empleando comunicación en tiempo real vía WebSockets/MQTT. La base de datos relacional (PostgreSQL) almacena la información de rutas, vehículos, usuarios y recorridos. Se realizaron pruebas funcionales y de rendimiento, incluyendo simulación de envíos periódicos de coordenadas GPS cada 5 segundos desde dispositivos móviles y la suscripción de clientes web a las actualizaciones en vivo. 
\textbf{Resultados:} UrbanTracker permite visualizar en un mapa la ubicación de los autobuses en servicio con una frecuencia de actualización de \textasciitilde5 segundos, superando el requisito mínimo de 10 segundos. La latencia observada entre el envío de una posición desde el móvil del conductor y su visualización en la aplicación de usuario es inferior a 2 segundos bajo condiciones de red normales. El sistema soporta las funcionalidades clave: consulta de rutas, vista de vehículos en tiempo real, inicio/cierre de recorrido por conductores, y administración completa de rutas y flotas. En pruebas de carga iniciales, la arquitectura basada en 
\textbf{publish/subscribe} demostró escalabilidad, en línea con reportes previos en sistemas basados en MQTT\cite{ortiz2022smolcomm}. 
\textbf{Discusión:} UrbanTracker logró implementar con éxito los requisitos definidos, brindando a los pasajeros información actualizada del servicio y a los administradores herramientas para la toma de decisiones informadas. Se discute la comparación con soluciones existentes y las implicaciones de la arquitectura tecnológica elegida, así como consideraciones de seguridad y privacidad en el manejo de datos de ubicación. 
\textbf{Conclusiones:} UrbanTracker evidencia que es factible mejorar la experiencia del transporte público local mediante geolocalización en tiempo real usando infraestructura asequible (smartphones y protocolos ligeros). Los resultados sugieren un impacto positivo en la reducción de tiempos de espera y en la eficiencia operativa. Trabajos futuros incluirán la integración con más fuentes de datos en tiempo real, refinamientos en la interfaz de usuario y evaluaciones de usabilidad a mayor escala.
